{\begingroup



\section{Categorical and Sheaf-Theoretic Preliminaries}

The coherent--constructible correspondence is an equivalence between
two derived categories of sheaves.  On the coherent side one considers
equivariant coherent sheaves on a toric variety, while on the
constructible side one considers constructible sheaves whose singular
support is constrained by a conic Lagrangian subset determined by the
fan.  Since morphisms in these categories are naturally derived
morphisms, it is not sufficient to regard them merely as ordinary
categories.  We will therefore use the language of categories enriched
over complexes, or equivalently dg categories, in order to keep track of
Hom-complexes, shifts, cones, and derived equivalences.

The purpose of this section is only to fix the categorical conventions
needed later.  We recall the notion of an enriched category, specialize
to dg categories, and explain how their homotopy categories recover the
triangulated categories appearing in the statement of the CCC.  We then
review the sheaf-theoretic categories used on both sides of the
correspondence: coherent and equivariant coherent sheaves, constructible
sheaves, standard and costandard sheaves, singular support, and sheaves
with prescribed singular support.

\subsection{dg Categories}

We first recall the notion of an enriched category.  In an ordinary
category, the morphisms between two objects form a set.  In many
geometric situations, however, the morphisms carry additional structure.
For example, in derived categories, morphisms are naturally organized
into complexes, whose cohomology groups give the corresponding Ext
groups.  This motivates the notion of a category enriched over a
monoidal category.

A \textbi{monoidal category} is a category \(\mathcal V\) together with
a bifunctor \(\otimes:\mathcal V\times\mathcal V\to\mathcal V\), a unit
object \(\mathbf 1\in\mathcal V\), and natural isomorphisms
\[
\alpha_{X,Y,Z}:(X\otimes Y)\otimes Z\xrightarrow{\sim}X\otimes(Y\otimes Z),
\qquad
\lambda_X:\mathbf 1\otimes X\xrightarrow{\sim}X,
\qquad
\rho_X:X\otimes\mathbf 1\xrightarrow{\sim}X,
\]
satisfying the pentagon and triangle axioms.  The pentagon axiom is the
commutativity of
\[
\begin{tikzcd}[column sep=large, row sep=large]
{((W\otimes X)\otimes Y)\otimes Z}
    \arrow[r, "\alpha_{W\otimes X,Y,Z}"]
    \arrow[d, "\alpha_{W,X,Y}\otimes \operatorname{id}_Z"']
&
{(W\otimes X)\otimes (Y\otimes Z)}
    \arrow[r, "\alpha_{W,X,Y\otimes Z}"]
&
{W\otimes (X\otimes (Y\otimes Z))}
\\
{(W\otimes (X\otimes Y))\otimes Z}
    \arrow[r, "\alpha_{W,X\otimes Y,Z}"']
&
{W\otimes ((X\otimes Y)\otimes Z)}
    \arrow[ru, "\operatorname{id}_W\otimes \alpha_{X,Y,Z}"']
&
\end{tikzcd}
\]
and the triangle axiom is the commutativity of
\[
\begin{tikzcd}[column sep=large, row sep=large]
{(X\otimes \mathbf 1)\otimes Y}
    \arrow[rr, "\alpha_{X,\mathbf 1,Y}"]
    \arrow[dr, "\rho_X\otimes \operatorname{id}_Y"']
&&
{X\otimes(\mathbf 1\otimes Y)}
    \arrow[dl, "\operatorname{id}_X\otimes \lambda_Y"]
\\
&
{X\otimes Y}.
&
\end{tikzcd}
\]
The bifunctor \(\otimes\) is called the \textbi{tensor product}.

Basic examples are \(\mathbf{Set}\) with the Cartesian product,
\(\mathbf{Ab}\) with the tensor product of abelian groups, and
\(\operatorname{Ch}(k)\) with the tensor product of complexes of
\(k\)-modules.  In \(\operatorname{Ch}(k)\), the unit object is the
complex \(k[0]\), namely \(k\) concentrated in degree zero.

A \textbi{\(\mathcal V\)-enriched category} is a category-like object
whose morphisms are objects of \(\mathcal V\), rather than sets.  More
precisely, a \(\mathcal V\)-enriched category \(\mathcal C\) consists of
a collection of objects \(\operatorname{Ob}(\mathcal C)\), an object
\(\operatorname{Hom}_{\mathcal C}(A,B)\in\mathcal V\) for every pair
\(A,B\in\mathcal C\), composition morphisms
\[
\operatorname{Hom}_{\mathcal C}(B,C)\otimes
\operatorname{Hom}_{\mathcal C}(A,B)
\longrightarrow
\operatorname{Hom}_{\mathcal C}(A,C),
\]
and unit morphisms \(\mathbf 1\to\operatorname{Hom}_{\mathcal C}(A,A)\).
These data are required to satisfy the usual associativity and unit
axioms, written in terms of the monoidal structure of \(\mathcal V\).

Ordinary categories are recovered by taking \(\mathcal V=\mathbf{Set}\)
with its Cartesian monoidal structure.  A \textbi{simplicial category}
is a category enriched in simplicial sets, so its morphisms are mapping
spaces.  A \textbi{dg category}, short for \textbi{differential graded
category}, over a field \(k\) is a category enriched in
\(\operatorname{Ch}(k)\), the monoidal category of cochain complexes of
\(k\)-vector spaces.

For a cochain complex \(V^\bullet\), its \(i\)-th cohomology is
\[
H^i(V^\bullet)
=
\frac{\ker(d^i:V^i\to V^{i+1})}
{\operatorname{im}(d^{i-1}:V^{i-1}\to V^i)}.
\]
A morphism of complexes \(f:V^\bullet\to W^\bullet\) is called a
\textbi{quasi-isomorphism} if the induced maps
\(H^i(f):H^i(V^\bullet)\to H^i(W^\bullet)\) are isomorphisms for all
\(i\in\mathbb Z\).

For a dg category \(\mathcal C\), its \textbi{homotopy category}
\(H^0(\mathcal C)\) is the ordinary category with the same objects as
\(\mathcal C\), and with morphism spaces
\[
\operatorname{Hom}_{H^0(\mathcal C)}(A,B)
=
H^0\operatorname{Hom}_{\mathcal C}^{\bullet}(A,B).
\]
A \textbi{dg functor} \(F:\mathcal C\to\mathcal D\) sends objects to
objects and induces morphisms of complexes
\[
\operatorname{Hom}_{\mathcal C}^{\bullet}(A,B)
\longrightarrow
\operatorname{Hom}_{\mathcal D}^{\bullet}(F(A),F(B)),
\]
compatible with composition and units.

A dg functor \(F:\mathcal C\to\mathcal D\) is called a
\textbi{quasi-equivalence} if, for all objects \(A,B\in\mathcal C\), the
induced morphism of Hom-complexes is a quasi-isomorphism, and if the
induced functor \(H^0(F):H^0(\mathcal C)\to H^0(\mathcal D)\) is
essentially surjective.  Equivalently, for every object \(D\in\mathcal D\),
there exists an object \(C\in\mathcal C\) such that \(F(C)\cong D\) in
\(H^0(\mathcal D)\).

A \textbi{right dg module} over \(\mathcal C\) is a dg functor
\(M:\mathcal C^{op}\to\operatorname{Ch}(k)\).  Equivalently, it assigns
to each object \(A\in\mathcal C\) a complex \(M(A)\), together with
action maps
\[
M(B)\otimes \operatorname{Hom}_{\mathcal C}^{\bullet}(A,B)
\longrightarrow M(A),
\]
compatible with composition and units.  We denote by
\(\operatorname{Mod}_{dg}(\mathcal C)\) the dg category of right dg
modules over \(\mathcal C\).

For each object \(A\in\mathcal C\), the associated \textbi{right
representable dg module} is
\(h_A:=\operatorname{Hom}_{\mathcal C}^{\bullet}(-,A)\).  Explicitly,
\(h_A(X)=\operatorname{Hom}_{\mathcal C}^{\bullet}(X,A)\), and a
morphism \(u:X\to Y\) acts by precomposition:
\[
\operatorname{Hom}_{\mathcal C}^{\bullet}(Y,A)
\longrightarrow
\operatorname{Hom}_{\mathcal C}^{\bullet}(X,A),
\qquad
g\longmapsto g\circ u.
\]

If \(V^\bullet\) is a cochain complex, its \textbi{shift} \(V[1]\) is
defined by \(V[1]^i=V^{i+1}\) and \(d_{V[1]}^i=-d_V^{i+1}\).  If
\(f:V^\bullet\to W^\bullet\) is a morphism of complexes, its
\textbi{cone} is the complex
\[
\operatorname{Cone}(f)^i=W^i\oplus V^{i+1},
\qquad
d(w,v)=(d_Ww+f(v),-d_Vv).
\]
The dg category \(\operatorname{Mod}_{dg}(\mathcal C)\) has shifts and
cones defined objectwise.  For example, \(M[1](A)=M(A)[1]\), and if
\(\varphi:M\to N\) is a closed degree-zero morphism of dg modules, then
\[
\operatorname{Cone}(\varphi)(A)
=
\operatorname{Cone}(M(A)\to N(A)).
\]

A dg category \(\mathcal C\) is called \textbi{pretriangulated} if
shifts and cones of representable dg modules are again representable.
More precisely, for every \(A\in\mathcal C\), the shifted module
\(h_A[1]\) is represented by an object of \(\mathcal C\), and for every
closed degree-zero morphism \(f:A\to B\), the cone of the induced
morphism \(h_A\to h_B\) is also represented by an object of
\(\mathcal C\).  In this case, \(H^0(\mathcal C)\) becomes a
triangulated category.

Not every dg category is pretriangulated.  Let \(\mathcal C\) be the dg
category with one object \(*\) and endomorphism complex
\(\operatorname{Hom}_{\mathcal C}(*,*)=k\), where \(k\) is concentrated
in degree zero.  The unique representable right dg module is
\(h_*=\operatorname{Hom}_{\mathcal C}(-,*)\simeq k\).  However, the
shifted module \(h_*[1]\simeq k[1]\) is not representable by an object
of \(\mathcal C\), since \(\mathcal C\) has only the object \(*\), whose
representable module is \(k\), not \(k[1]\).  Therefore \(\mathcal C\)
is not pretriangulated.

The \textbi{dg Yoneda embedding} is the dg functor
\[
Y:\mathcal C\longrightarrow \operatorname{Mod}_{dg}(\mathcal C),
\qquad
A\longmapsto \operatorname{Hom}_{\mathcal C}^{\bullet}(-,A).
\]
The \textbi{pretriangulated envelope} of \(\mathcal C\), denoted
\(\operatorname{Pretr}(\mathcal C)\), is the smallest full dg subcategory
of \(\operatorname{Mod}_{dg}(\mathcal C)\) which contains all
representable dg modules \(h_A\) and is closed under finite direct sums,
shifts, and cones.

For the one-object dg category above, the pretriangulated envelope is
obtained by adjoining formal shifts, cones, and finite direct sums.
After idempotent completion, it is identified with
\(\operatorname{Perf}(k)\), the dg category of perfect complexes over
\(k\).

Let \(\mathcal C\) be a pretriangulated dg category.  A full dg
subcategory \(\mathcal A\subset\mathcal C\) is called \textbi{thick} if
\(H^0(\mathcal A)\) is a full triangulated subcategory of
\(H^0(\mathcal C)\) closed under direct summands.  Equivalently, if
\(X\in H^0(\mathcal A)\) and \(Y\) is a retract of \(X\) in
\(H^0(\mathcal C)\), then \(Y\in H^0(\mathcal A)\).

For a collection of objects \(\mathcal S\subset\operatorname{Ob}(\mathcal C)\),
we denote by \(\langle\mathcal S\rangle\) the smallest full dg subcategory
of \(\mathcal C\) whose homotopy category is a thick triangulated
subcategory of \(H^0(\mathcal C)\) containing \(\mathcal S\).

Let \(\mathcal C\) and \(\mathcal D\) be pretriangulated dg categories.
A dg functor \(F:\mathcal C\to\mathcal D\) is called \textbi{exact} if
the induced functor \(H^0(F):H^0(\mathcal C)\to H^0(\mathcal D)\) is an
exact functor of triangulated categories.  Equivalently, it commutes
with shifts up to isomorphism and sends distinguished triangles to
distinguished triangles.

In practice, we will use the following criterion.

\begin{theorem}[Generator Criterion for Quasi-Equivalences]\label{thm:generator-criterion}
Let \(F:\mathcal C\to\mathcal D\) be an exact dg functor between
pretriangulated dg categories.  Let
\(\mathcal S\subset\operatorname{Ob}(\mathcal C)\) be a collection of
objects such that \(\langle\mathcal S\rangle=\mathcal C\) and
\(\langle F(\mathcal S)\rangle=\mathcal D\).  Suppose that, for all
\(A,B\in\mathcal S\), the induced morphism of complexes
\[
\operatorname{Hom}_{\mathcal C}^{\bullet}(A,B)
\longrightarrow
\operatorname{Hom}_{\mathcal D}^{\bullet}(F(A),F(B))
\]
is a quasi-isomorphism.  Then \(F\) is a quasi-equivalence.
\end{theorem}

Therefore, the proof of a dg quasi-equivalence can often be reduced to
two tasks:
\begin{enumerate}
    \item compare the Hom-complexes on a set of generators;
    \item prove that these generators generate both sides.
\end{enumerate}





\subsection{Coherent and Equivariant Coherent Sheaves}

Let \(X\) be a locally Noetherian scheme.  We denote by
\(\operatorname{Coh}(X)\) the abelian category of coherent
\(\mathcal O_X\)-modules.  Let
\(\operatorname{Ch}^b(\operatorname{Coh}(X))\) be the category of
bounded complexes of coherent sheaves, and let
\(K^b(\operatorname{Coh}(X))\) be its homotopy category.  Thus
\(K^b(\operatorname{Coh}(X))\) has the same objects as
\(\operatorname{Ch}^b(\operatorname{Coh}(X))\), but its morphisms are
chain maps modulo chain homotopy.

The \textbi{bounded derived category} of coherent sheaves is the Verdier
localization
\[
D^b(\operatorname{Coh}(X))
=
K^b(\operatorname{Coh}(X))[\mathcal W^{-1}],
\]
where \(\mathcal W\) denotes the class of quasi-isomorphisms in
\(K^b(\operatorname{Coh}(X))\).

Now suppose that an algebraic group \(G\) acts on \(X\).  Let
\(a:G\times X\to X\) be the action map, \(p_X:G\times X\to X\) the
projection to \(X\), and \(\mu:G\times G\to G\) the multiplication map.
On \(G\times G\times X\), define
\[
q_X(g,h,x)=x,\qquad
q_{23}(g,h,x)=(h,x),
\qquad
(\operatorname{id}_G\times a)(g,h,x)=(g,hx),
\]
and regard \(\mu\times \operatorname{id}_X:G\times G\times X\to G\times X\)
as the map \((g,h,x)\mapsto (gh,x)\).

A \textbi{\(G\)-equivariant coherent sheaf} on \(X\) is a coherent sheaf
\(\mathcal F\) together with an isomorphism
\[
\phi:p_X^*\mathcal F\xrightarrow{\sim}a^*\mathcal F
\]
satisfying the following cocycle condition:
\[
\begin{tikzcd}[column sep=large, row sep=large]
q_X^*\mathcal F
    \arrow[r, "q_{23}^*\phi"]
    \arrow[dr, "{(\mu\times \operatorname{id}_X)^*\phi}"']
&
(a\circ q_{23})^*\mathcal F
    \arrow[d, "{(\operatorname{id}_G\times a)^*\phi}"]
\\
&
(a\circ(\operatorname{id}_G\times a))^*\mathcal F .
\end{tikzcd}
\]
Here we use the equality
\(a\circ(\operatorname{id}_G\times a)=a\circ(\mu\times \operatorname{id}_X)\)
as maps \(G\times G\times X\to X\). We call such an isomorphism \(\phi\) a \textbi{\(G\)-linearization of \(\mathcal F\)}.

We denote by \(\operatorname{Coh}_G(X)\) the \textbi{abelian category of
\(G\)-equivariant coherent sheaves on \(X\)}, and by
\(D^b(\operatorname{Coh}_G(X))\) its \textbi{bounded equivariant derived
category}.

A complex \(\mathcal E^\bullet\) of \(\mathcal O_X\)-modules is called
\textbi{perfect} if it is locally quasi-isomorphic to a bounded complex
of finite rank locally free sheaves.  We write \(\operatorname{Perf}(X)\)
for the \textbi{dg category of perfect complexes on \(X\)}.  If \(G\)
acts on \(X\), we similarly write \(\operatorname{Perf}_G(X)\) for the
\textbi{dg category of \(G\)-equivariant perfect complexes}.

The following standard fact explains why perfect complexes provide the
correct dg model for the coherent side of the correspondence.  In the
smooth case, bounded coherent complexes can be represented by perfect
complexes.  In the equivariant setting, the same statement holds
provided that there are sufficiently many equivariant vector bundles.
\begin{theorem}
\label{thm:perfect-complexes-as-compact-objects}
Let \(X\) be a quasi-compact and quasi-separated scheme. Then
\[
D(\operatorname{QCoh}(X))^c \simeq \operatorname{Perf}(X).
\]
Equivariantly, for the toric varieties considered in this note, one has
\[
D(\operatorname{QCoh}_T(X_\Sigma))^c
\simeq
\operatorname{Perf}_T(X_\Sigma).
\]
\end{theorem}
\begin{theorem}
Let \(X\) be a smooth Noetherian scheme over a field \(k\).
\begin{enumerate}
    \item The dg category \(\operatorname{Perf}(X)\) is a dg enhancement
    of \(D^b(\operatorname{Coh}(X))\).  Equivalently,
    \[
    H^0(\operatorname{Perf}(X))
    \simeq
    D^b(\operatorname{Coh}(X)).
    \]

    \item Let \(G\) be an algebraic group acting on \(X\).  If \(X\) has
    the equivariant resolution property, then
    \(\operatorname{Perf}_G(X)\) is a dg enhancement of
    \(D^b(\operatorname{Coh}_G(X))\).  Equivalently,
    \[
    H^0(\operatorname{Perf}_G(X))
    \simeq
    D^b(\operatorname{Coh}_G(X)).
    \]
\end{enumerate}
\end{theorem}

Here, we say that \(X\) has the \textbi{equivariant resolution property}
if every \(G\)-equivariant coherent sheaf
\(\mathcal F\in \operatorname{Coh}_G(X)\) admits a surjection
\[
\mathcal E\twoheadrightarrow \mathcal F
\]
from a \(G\)-equivariant vector bundle \(\mathcal E\) on \(X\).



\subsection{Constructible Sheaves}

Let \(M\) be a real analytic manifold, and let \(k\) be a fixed field of
coefficients.  We denote by \(Sh(M)\) the \textbi{dg category of
complexes of sheaves of \(k\)-vector spaces on \(M\)}.  Equivalently,
\(Sh(M)\) may be regarded as a dg enhancement of the derived category of
sheaves of \(k\)-vector spaces on \(M\).

We first recall the notion of a locally constant sheaf.  A sheaf
\(\mathcal F\) of \(k\)-vector spaces on \(M\) is called
\textbi{locally constant} if every point \(x\in M\) has an open
neighborhood \(U\) such that \(\mathcal F|_U\simeq k_U^{\oplus r}\) for
some integer \(r\) depending on \(U\).  A \textbi{local system} is a
locally constant sheaf whose stalks are finite-dimensional
\(k\)-vector spaces.

A \textbi{stratification} of \(M\) is a locally finite decomposition
\(M=\bigsqcup_{\alpha\in A}S_\alpha\) into locally closed smooth
submanifolds.  A sheaf \(\mathcal F\) is \textbi{constructible} with
respect to this stratification if each restriction
\(\mathcal F|_{S_\alpha}\) is a local system.  A sheaf is called
constructible if it is constructible with respect to some such
stratification.

We denote by \(Sh_c(M)\) the full dg subcategory of \(Sh(M)\) consisting
of constructible complexes.  More precisely, a complex
\(\mathcal F^\bullet\in Sh(M)\) is constructible if each cohomology
sheaf \(\mathcal H^i(\mathcal F^\bullet)\) is constructible and
\(\mathcal H^i(\mathcal F^\bullet)=0\) for all but finitely many
\(i\).  We write \(Sh_{cc}(M)\) for the full dg subcategory of
\(Sh_c(M)\) consisting of constructible complexes with compact support.
Here the \textbi{support of a complex} \(\mathcal F^\bullet\) is the
union of the supports of its cohomology sheaves.

For a fixed stratification \(\mathcal S=\{S_\alpha\}_{\alpha\in A}\), we
write \(D^b_{\mathcal S}(M;k)\) for the bounded derived category of
\(k\)-sheaves constructible with respect to \(\mathcal S\).




Let \(j:U\hookrightarrow M\) be the inclusion of a locally closed subset.
We denote by \(k_U\) the constant sheaf on \(U\) with coefficients in
\(k\).  The \textbi{standard sheaf associated to \(U\)} is \(j_!k_U\),
where \(j_!\) denotes extension by zero.  The \textbi{costandard sheaf
associated to \(U\)} is \(j_*k_U\), where \(j_*\) denotes the usual
direct image.

Standard and costandard sheaves are the basic building blocks of
constructible categories.  The following generation statement is the
constructible-side analogue of choosing generators on the coherent side.


\begin{theorem}[Generation by standard and costandard sheaves]
Let \(M\) be a real analytic manifold, and let
\(M=\bigsqcup_{\alpha\in A}S_\alpha\) be a finite stratification by
locally closed submanifolds.  Let \(j_\alpha:S_\alpha\hookrightarrow M\)
denote the inclusion.  Then
\[
D^b_{\mathcal S}(M;k)
=
\left\langle
j_{\alpha !}L
\;\middle|\;
\alpha\in A,\ 
L \text{ a finite-dimensional local system on }S_\alpha
\right\rangle.
\]
Equivalently,
\[
D^b_{\mathcal S}(M;k)
=
\left\langle
j_{\alpha *}L
\;\middle|\;
\alpha\in A,\ 
L \text{ a finite-dimensional local system on }S_\alpha
\right\rangle.
\]
In particular, if every stratum \(S_\alpha\) is contractible, then
\[
D^b_{\mathcal S}(M;k)
=
\left\langle
j_{\alpha !}k_{S_\alpha}
\;\middle|\;
\alpha\in A
\right\rangle
=
\left\langle
j_{\alpha *}k_{S_\alpha}
\;\middle|\;
\alpha\in A
\right\rangle.
\]
\end{theorem}
\subsection{Singular Support}

Let \(M\) be a real analytic manifold, and let \(T^*M\) denote its
cotangent bundle.  The cotangent bundle carries a natural fiberwise
dilation action of \(\mathbb R_{>0}\), given by
\(t\cdot(x,\xi)=(x,t\xi)\).  A subset
\(\Lambda\subset T^*M\) is called \textbi{conic} if it is invariant
under this action; equivalently, \((x,\xi)\in\Lambda\) implies
\((x,t\xi)\in\Lambda\) for every \(t\in\mathbb R_{>0}\).

For a constructible complex of sheaves \(\mathcal F\in Sh_c(M)\), its
\textbi{singular support}, also called its \textbi{microsupport}, is a
closed conic subset \(SS(\mathcal F)\subset T^*M\).  It is a microlocal
refinement of the ordinary support.  A covector
\((x_0,\xi_0)\in T^*M\) is not contained in \(SS(\mathcal F)\) if there
exists a conic open neighborhood \(U\subset T^*M\) of \((x_0,\xi_0)\)
such that, for every \(C^1\)-function \(\varphi\) defined near \(x_0\)
with \(\varphi(x_0)=0\) and \((x_0,d\varphi_{x_0})\in U\), one has
\[
\bigl(R\Gamma_{\{\varphi\geq 0\}}(\mathcal F)\bigr)_{x_0}=0.
\]
Here \(R\Gamma_Z(\mathcal F)\) denotes the derived functor of sections
with support in a closed subset \(Z\subset M\).  Thus
\(\bigl(R\Gamma_{\{\varphi\geq 0\}}(\mathcal F)\bigr)_{x_0}\) is the
local cohomology complex of \(\mathcal F\) at \(x_0\) with support in the
half-space \(\{\varphi\geq 0\}\).

Equivalently, \(SS(\mathcal F)\) is the closed conic subset of \(T^*M\)
consisting of covectors along which this local cohomological vanishing
fails.  It records the cotangent directions in which sections of
\(\mathcal F\) cannot be propagated microlocally.

If \(\mathcal F\) is locally constant on \(M\), then it has no nonzero
singular directions; hence \(SS(\mathcal F)\subset T_M^*M\), where
\(T_M^*M\subset T^*M\) denotes the zero section.

Let \(S\subset M\) be a locally closed submanifold.  The
\textbi{conormal bundle} of \(S\) in \(M\) is
\[
T_S^*M
=
\{(x,\xi)\in T^*M\mid x\in S,\ \xi|_{T_xS}=0\}.
\]
If \(\mathcal F\) is constructible with respect to a stratification
\(M=\bigsqcup_{\alpha\in A}S_\alpha\), then
\[
SS(\mathcal F)\subset \bigcup_{\alpha\in A}T^*_{S_\alpha}M.
\]
This gives the basic geometric control on the singularities of a
constructible sheaf.

Let \(\Lambda\subset T^*M\) be a closed conic subset.  We define
\[
Sh_\Lambda(M)
:=
\{\mathcal F\in Sh_c(M)\mid SS(\mathcal F)\subset \Lambda\}
\]
to be the full dg subcategory of \(Sh_c(M)\) consisting of constructible
complexes whose singular support is contained in \(\Lambda\).  Thus
\(\Lambda\) imposes a microlocal constraint on constructible sheaves:
instead of allowing arbitrary constructible singularities, one only
allows singular behavior in the cotangent directions prescribed by
\(\Lambda\).

\begin{theorem}
Let \(\Lambda\subset T^*M\) be a closed conic subset.  The full
subcategory \(Sh_\Lambda(M)\subset Sh_c(M)\) is closed under shifts,
cones, finite direct sums, and retracts.  Consequently,
\(H^0(Sh_\Lambda(M))\) is a triangulated subcategory of
\(H^0(Sh_c(M))\).
\end{theorem}
\begin{theorem}[Microlocal Morse lemma]
Let \(M\) be a real analytic manifold, let
\(F\in \operatorname{Sh}_c(M)\), and let \(f:M\to \mathbb R\) be a real
analytic function. Let \(a<b\). Suppose that \(f\) is proper on
\(\operatorname{supp}(F)\cap f^{-1}([a,b])\), and that
\[
df(x)\notin SS(F)
\]
for every \(x\in \operatorname{supp}(F)\cap f^{-1}([a,b])\). Then the
restriction map
\[
R\Gamma(f^{-1}((-\infty,b));F)
\longrightarrow
R\Gamma(f^{-1}((-\infty,a));F)
\]
is a quasi-isomorphism.
\end{theorem}

Thus singular support controls the cotangent directions in which local
cohomology may fail to propagate. In directions outside \(SS(F)\), passing
through nearby level sets of \(f\) does not create new cohomology.
\endgroup}